\section{A bound for the height error at each member of a window}
\label{sec:height-gap}

The finite Fourier calculation uses an upper bound for a difference
\(\mathrm{ASK}-\mathrm{Main}\). It is important that this bound holds at
each member of a finite family. An upper bound on its mean over the
family would leave individual prime sums uncontrolled. We translate
the pointwise slope calculation which produces the error before its
terms are estimated by powers of the scale.

Fix a member of the family and write \(W\ne\varnothing\) for its
finite set of words and \(d_u>0\) for the weight of \(u\in W\).
Write \(\langle f\rangle=\sum d_u f_u/\sum d_u\). For each word let
\(\nu_u,\nu_u^+,\mu_u,\mu_u^+\) denote respectively the slope and
maximal slope of its target and source lattices. Write \(h_u\) for its
height, \(T_u>0\) for its count, and \(E_u\) for its exponent sum.
Let \(G,F,\mu_C,\mu_C^{\min}\) be the quantities denoted
\(\mathrm{Ga},\mathrm{Fo},\mu_{\mathrm{CC}},\mu_{\min,\mathrm{CC}}\)
at this member. Set
\[
 s=\langle\nu_u-h_u\rangle,\qquad
 z=\max_{u\in W}(\nu_u^+-h_u)-s .
\]
Since \(\nu_u\le\nu_u^+\), we have \(z\ge0\).

Let there be \(n_T\) top indices \(t\) with weights \(a_t\ge0\).
Put \(A=\sum_t a_t\). Write \(m\) for the chain length, \(f>0\) for
the number of geometric components, \(v>0\) for the volume parameter,
\(w_{j\rho}>0\) for the local weights, and \(g\) for the exponent
parameter \(\mathrm{lg}\). The notation \(L_a,L,a,b,b',k\) and
\(B_0,N_c,X''\) is local to this section, with \(k,b>0\). The
constants denoted by \(c\) with subscripts below are those of the
window in the source. We keep separate the upper error in the count,
the cost of division, and the loss in passing to the quotient. Put
\begin{align*}
 C_T&=c_{\rm count}(L_a+m^2L_a/k),\\
 \lambda&=\sum_{j=1}^m f^{-1}\sum_{\rho=1}^f
            \log(w_{j\rho}v/\pi)+mg-C_T,\\
 E_D&=c_{\rm div}n_TaL_a+
          c_{\rm div}m(B_0+N_c)x^{e_{\rm div}\varepsilon},\\
 S_-&=c_{\rm slope}n_TaL_a(1+m/k),\qquad
 S_+=c_{\rm slope}n_TaL_a,\\
 P&=c_{\rm loss}n_Tmbx^{e_{\rm loss}\varepsilon},\\
 L_u&=c_1\mathrm{IO}+c_1 n_Tm(a/k+b')\log T_u+c_1 X''+1 .
\end{align*}
Here \(\mathrm{IO}\) is the local loss and \(c_1,n_T,m,a,b'\ge0\).
The two quotient-slope comparisons read
\begin{align*}
 \nu^+_u&\le\mu_u^++G+F-E_u+c_1mB_0,\\
 \nu_u&\ge\mu_u+G+F-E_u-L_u .
\end{align*}
The slope bounds give \(-S_-\le\mu_u\) and \(\mu^+_u\le S_+\).
The exponent, height and count comparisons used here give
\begin{align*}
 |E_u-mAg|&\le c_E m B_0,\\
 |(h_u+A\log T_u)-(h_{u'}+A\log T_{u'})|
       &\le c_H(n_T a L_a+mB_0),\\
 \lambda\le\log T_u&\le c_{\rm count}mL,\qquad
 |\log T_u-\log T_{u'}|\le c_{\rm count}\log k.
\end{align*}
Assume also \(\mathrm{ASK}-G\le P,\ \mathrm{IO}\le P\). These
pointwise hypotheses are quantitative inputs from earlier parts of
the construction.

\paragraph{Cancellation of the common exponent.}
Put \(\omega=\gamma+\log(\pi/v)-\mathcal L\), where
\(\mathcal L\) denotes the source's logarithmic-derivative parameter.
Let
\[
 D_0=m\sum_t(a_t\log a_t-a_t-a_t\omega).
\]
Terms with \(a_t=0\) are interpreted by continuity. The division
inequality is \(\mu_C+h_u-F\le D_0-A\log T_u+E_D\).
Insert it in the lower quotient comparison. For every \(u\) this
gives
\[
 \nu_u-h_u\ge G+\mu_C
          -S_- -m(Ag+c_EB_0)-D_0+A\lambda-E_D-L_u .
\]
Expanding \(\lambda\) and \(\omega\) yields the identity
\begin{align*}
 D_0-A\lambda+mAg
  &=\mathrm{Main}+AC_T,\\
 \mathrm{Main}
  &=\sum_{j=1}^m\sum_t a_t\!
  \left(\log a_t-1-\gamma+\mathcal L
               - f^{-1}\sum_\rho\log w_{j\rho}\right).
\end{align*}
Indeed the \(mg\) terms cancel first, and
\(\log(\pi/v)+\log(w_{j\rho}v/\pi)=\log w_{j\rho}\).
Taking the weighted mean proves
\begin{equation}
  G+\mu_C-\mathrm{Main}-K_0-\langle L\rangle \le s,\qquad
  K_0=AC_T+E_D+mc_EB_0+S_- .
  \label{eq:height-gap-lower}
\end{equation}

\paragraph{Variation from word to word.}
For two words \(u,u'\), subtracting the two exponent bounds and the
height comparisons gives
\[
  (-E_u-h_u)-(-E_{u'}-h_{u'})\le V,\quad
  V=2mc_EB_0+c_H(n_TaL_a+mB_0)+Ac_{\rm count}\log k.
\]
Subtract the lower quotient comparison at \(u'\) from the upper
one at \(u\). The terms \(G,F\) cancel. After subtracting the heights
the result is at most \(K_3+L_{u'}\), with
\(K_3=S_-+S_++V+c_1mB_0\). Taking the mean over \(u'\) and the
maximum over \(u\) yields \(z\le K_3+\langle L\rangle\). By the
upper bound on \(\log T_u\) and the bound on \(\mathrm{IO}\),
\[
 \langle L\rangle\le M_L
    :=c_1P+c_1 n_Tm(a/k+b')c_{\rm count}mL+c_1X''+1 .
\]

\paragraph{The logarithmic term.}
The slope inequality on the common lattice gives
\[
 s-\mu_C\le\log |W|+
 \frac{4z\log(1+H/z)+(\eta+4\sqrt e)H}{1-\delta}
\]
when \(s-\mu_C^{\min}\le H,\ H>0,\ 0\le\delta<1,\ 0\le e\le1/4\)
and \(\eta\ge0\). At \(z=0\) use the continuous value zero for the
product. Suppose \(H\le Xb\), \(X=x^{c_{HH}}>0\).
For \(z>0\), the tangent-line inequality for the logarithm gives
\[
 z\log(1+H/z)
 \le z\log(1+X)+\frac{H-Xz}{1+X}
 \le (z+b)\log(1+X).
\]
The last inequality uses \(X/(1+X)\le\log(1+X)\), by integrating
\((1+t)^{-1}\) from \(0\) to \(X\). This holds at \(z=0\) as well.
If the window gives uniform bounds
\(\delta\le\bar\delta<1,\ \eta\le\bar\eta,\ e\le\bar e\le1/4\),
the nonnegativity of all numerator terms gives
\begin{align}
 \mathrm{ASK}-\mathrm{Main}
  &\le K_0+\log|W|+P+M_L \notag\\
  &\quad+
    \frac{4\log(1+X)(K_3+M_L+b)
                     +(\bar\eta+4\sqrt{\bar e})Xb}{1-\bar\delta}.
    \label{eq:unit-gap}
\end{align}
This follows by adding \eqref{eq:height-gap-lower},
the bound for \(s-\mu_C\), and \(\mathrm{ASK}-G\le P\) with the
indicated signs. The right-hand side is the error denoted
\(\mathrm{ErrX5at}\) for this member. In particular \(\log|W|=0\)
is permitted. No averaging over different windows, scales or
members occurs in this derivation.
\footnote{\href{https://github.com/Deicyde/riemann-hypothesis/blob/40e56bab/src/Formalism/RankGapBounds/MeanBoundSecondForm.lean}{Slope, division and height hypotheses};
\href{https://github.com/Deicyde/riemann-hypothesis/blob/40e56bab/src/Formalism/SizeWindowBrackets/QuotientSlopeBoundsWith.lean}{count and quotient inequalities and the error formula};
\href{https://github.com/Deicyde/riemann-hypothesis/blob/40e56bab/src/Formalism/SizeWindowBrackets/ChainSlopeWindows/LevelOfCongruence.lean}{the pointwise deduction}.}
